\section{Diseño experimental y reproducibilidad}

\subsection{Escenarios de muestreo repetido}

La configuración fija
\begin{equation}
 S_0=100,
 \qquad \mu_0=0.08,
 \qquad r=0.03,
 \qquad q=0,
 \qquad \Delta t=1/252.
\end{equation}
Se consideran tres cantidades de información histórica,
\begin{equation}
 n\in\{63,252,1260\},
\end{equation}
que corresponden aproximadamente a 0.25, 1 y 5 años de retornos diarios; tres volatilidades verdaderas,
\begin{equation}
 \sigma_0\in\{0.15,0.25,0.40\};
\end{equation}
tres niveles de moneyness,
\begin{equation}
 K/S_0\in\{0.8,1.0,1.2\};
\end{equation}
y tres madureces,
\begin{equation}
 T\in\{0.5,1,2\}.
\end{equation}
El producto cartesiano genera 81 escenarios de evaluación. Para cada combinación $(n,\sigma_0)$ se simulan 2,000 historias independientes, dando lugar a 18,000 ajustes posteriores. Cada posterior se reutiliza sobre las nueve combinaciones $(K/S_0,T)$; no se repite la inferencia al cambiar términos contractuales.

Cada ajuste utiliza 20,000 iteraciones de Metropolis--Hastings, 4,000 de burn-in y thinning igual a 2. La semilla raíz es 20260909 y las semillas de unidades independientes se generan de forma determinista, de modo que una corrida puede reconstruirse a partir de la configuración y el commit de código.

\subsection{Benchmark neutral al riesgo}

Para cada contrato se precomputa una malla
\begin{equation}
 \sigma\in[0.02,1.00]
\end{equation}
con 161 puntos. Cada punto se valúa con dos millones de trayectorias bajo $\Qmeasure$, muestreo antitético y control geométrico. Los nueve grids se calcularon con GPU mediante CuPy. El error estándar Monte Carlo máximo reportado entre las nueve configuraciones contractuales se encuentra entre 0.00218 y 0.01805, con mediana de 0.00579 entre los máximos por contrato.

El benchmark de cada escenario es
\begin{equation}
 C_0=C^{\Qmeasure}(\sigma_0),
 \label{eq:trueprice}
\end{equation}
aproximado mediante interpolación en la malla de alta precisión. Por construcción, $C_0$ depende de la volatilidad verdadera y no de una posterior estimada. Los mismos grids se utilizan para evaluar \eqref{eq:fb}--\eqref{eq:mle}, reduciendo ruido numérico en comparaciones entre reglas.

\subsection{Métricas}

Para cada método $M$ y escenario, con $R=2000$ réplicas, se calculan
\begin{align}
 \Bias_M
 &=\frac1R\sum_{j=1}^{R}(\widehat C_{M,j}-C_0),\\
 \MAE_M
 &=\frac1R\sum_{j=1}^{R}|\widehat C_{M,j}-C_0|,\\
 \RMSE_M
 &=\sqrt{\frac1R\sum_{j=1}^{R}(\widehat C_{M,j}-C_0)^2}.
\end{align}
Para la posterior de precios se registra además si el intervalo central al 95\% contiene el benchmark. Las métricas se agregan por escenario y también se estudia directamente la diferencia
\begin{equation}
 \Delta_{FB,PM}=\widehat C_{FB}-\widehat C_{PM}.
 \label{eq:gap}
\end{equation}

\subsection{Aplicación empírica WTI}

La evidencia simulada se complementa con una aplicación a futuros de petróleo WTI. Se descarga el histórico diario de \texttt{CL=F} desde Yahoo Finance mediante \texttt{yfinance}. El símbolo se utiliza como proxy continuo/front-month y no como una reconstrucción propia de first-nearby, dado que la regla histórica de roll del proveedor no se observa de forma contractual en la descarga.

El snapshot completo contiene 6,539 registros entre 2000-08-23 y 2026-09-09. Debido a que el GBM requiere una variable de estado positiva y el mercado WTI presentó precios no positivos en 2020, la muestra de estimación se fija ex ante desde 2021-01-01. La primera observación efectiva es 2021-01-04 y la última 2026-09-09, con 1,429 precios y 1,428 log-retornos. Se utiliza nuevamente $\Delta t=1/252$.

Para esta aplicación se ejecutan 50,000 iteraciones de Metropolis--Hastings, con 10,000 de burn-in y 40,000 draws posteriores. El snapshot local se conserva fuera del repositorio, mientras el manifiesto versionado registra proveedor, fechas, versión de \texttt{yfinance}, tamaño muestral y SHA-256 del archivo. Esta parte del estudio se interpreta como aplicación empírica de la inferencia paramétrica, no como validación de precios de opciones asiáticas observadas.

\subsection{Infraestructura computacional}

La corrida principal corresponde al fingerprint reproducible \texttt{f497062f0c571126}. Se ejecutó con 24 procesos de CPU para las inferencias independientes y una NVIDIA RTX 4500 Ada Generation para las mallas de pricing mediante CuPy 14.2.0. El manifiesto conserva configuración, semilla, versiones de software, hardware y commit de Git. Los checkpoints se almacenan por unidad de inferencia y por punto de pricing, permitiendo reanudar una corrida interrumpida sin recalcular unidades completas. Los resultados agregados y el código de análisis permanecen versionados en el repositorio.

La aplicación WTI utiliza el mismo entorno reproducible y la semilla 20260909 para la cadena MCMC. El snapshot empleado se identifica por el SHA-256 \texttt{b55c92a4f354fa0f6d5a4c72907920f46a5cf6aebf0a6db5fc4a93d88dcb469d}, permitiendo distinguir de forma inequívoca los datos que produjeron los resultados empíricos aun si el proveedor actualiza posteriormente su histórico.
